\begin{textsample}{POS Dim 1 – pe – Score 34.00 – pe/pe\_em\_12.txt}  \label{ex:f1_pos_184}
1.9 TEMAS TRANSVERSAIS E INTEGRADORES DO CURRÍCULO O Currículo de Pernambuco contempla temas sociais e saberes \textbf{que} envolvem várias dimensões, como : política, social, histórica, cultural, ética e econômica.
Tais dimensões são necessárias à formação integral dos estudantes e afetam a vida humana em escala local, regional e global, trazendo temáticas \textbf{que} devem integrar o cotidiano da escola.
Alguns desses temas estão diretamente relacionados às legislações específicas, \textbf{enquanto} outros são sugeridos em diretrizes curriculares, ou mesmo, \textbf{demandados} pela própria comunidade educativa.
O \textbf{que} os une é o fato de se relacionarem a diferentes componentes curriculares, garantindo \textbf{uma} \textbf{abordagem} interdisciplinar, transversal e integradora.
Citamos alguns desses temas, entendendo \textbf{que} outros poderão ser acrescentados em função de novas demandas legais ou por escolha das próprias escolas, inserindo -os em seus projetos político-pedagógicos por meio de práticas educativas voltadas para a criação de \textbf{uma} cultura de paz.
A Educação em Direitos Humanos-EDH ( Plano Nacional de Educação em Direitos Humanos, 2006, Decreto nº 7.037/2009, Parecer CNE/CP nº 8/2012 e Resolução CNE/CP nº 1/2012 ), alicerçada no respeito e proteção à \textbf{dignidade} do ser humano, compreende o conjunto de práticas educativas fundamentadas nos direitos humanos, tendo como objetivo formar o \textbf{sujeito} de direito.
Nesse contexto, a Secretaria de Educação e Esportes de Pernambuco, nas últimas décadas, assumiu a EDH como \textbf{norteadora} das políticas educacionais do estado de Pernambuco e pautou -a no \textbf{compromisso} pela construção de \textbf{uma} escola \textbf{que} se reconheça como espaço pleno de vivências de direitos, premissa fundamental para \textbf{embasar} as relações humanas \textbf{que} acontecem na escola em todos os seus âmbitos.
As Diretrizes Nacionais para a Educação em Direitos Humanos ( Resolução CNE/CP nº 1/2012 ) prescrevem \textbf{que}, na Educação Básica, o currículo poderá ser estruturado tomando por base a perspectiva disciplinar, transversal ou mista, fundindo disciplinaridade e transversalidade.
Ao fazer a opção por tratar a EDH na perspectiva transversal, o estado de Pernambuco filia -se ao entendimento de \textbf{que} a cultura dos direitos humanos, conteúdo da EDH, não cabe apenas em \textbf{um} componente curricular, devendo, assim, ganhar espaço no conjunto dos componentes \textbf{que} compõem o currículo.
Materializada na perspectiva transversal, a EDH fortalece os paradigmas da educação integral, considerando os estudantes em todas as suas dimensões.
Além disso, sedimenta \textbf{uma} cultura de paz na escola, fundamentada na defesa e reconhecimento da igualdade de direitos, valorização das diferenças e das diversidades, laicidade do estado e democracia na educação.
A escola, na perspectiva da EDH, deve desenvolver \textbf{uma} educação pautada em várias dimensões necessárias à formação cidadã : ciências, artes, cultura, história, ética, afetividade, entre outras.
Assim, a escola é concebida como espaço sociocultural, lugar de convivência inclusiva, respeitosa e afetiva.
O ambiente escolar deve proporcionar, também, \textbf{uma} convivência acolhedora, de autorresponsabilidade com o desempenho de cada estudante, de cada professor, \textbf{consigo} mesmo, bem como de cuidado com o outro, considerando a \textbf{dignidade} de todo ser humano.
No campo da discussão dos Direitos da Criança e do Adolescente ( Lei nº 8.069/1990 – Estatuto da Criança e do Adolescente, Lei nº 12.852/2013 – Estatuto da Juventude, Lei nº 13.257/2016 – Marco Legal da Primeira Infância, de 08 de março de 2016 ), o direito de brincar da criança e também o direito de ser cuidada por profissionais \textbf{qualificados}, na primeira infância, devem ser prioridade nas políticas públicas.
A criança tem, sobretudo, o direito a ter a presença da mãe, pai e/ou cuidador em casa nos primeiros meses por meio da licença-maternidade e paternidade concedida para cumprimento dos cuidados.
Por sua vez, o direito à educação deve ser garantido a todas as crianças e adolescentes, observando o pleno desenvolvimento de suas potencialidades por meio de \textbf{uma} preparação cultural \textbf{qualificada}, \textbf{uma} base científica e humana na perspectiva de contribuir para a superação das desvantagens decorrentes das condições socioeconômicas e culturais adversas.
Nessa \textbf{direção}, situamos também o Estatuto da Juventude, \textbf{que} vem corroborar a inserção social \textbf{qualificada} do jovem como lei complementar ao Estatuto da Criança e do Adolescente, visando garantir direitos de quem tem entre 15 e 29 anos de idade.
O Estatuto da Juventude propõe expansão das garantias dadas à infância e à adolescência, além da compreensão de \textbf{que} o jovem deve ser \textbf{visto} nas suas necessidades no momento presente e não a posteriori.
Desse modo, as aprendizagens essenciais devem ser contempladas, proporcionando o desenvolvimento das competências e habilidades necessárias, e possibilitando às crianças, adolescentes e jovens o direito a \textbf{uma} educação de qualidade para \textbf{que} possam atuar socialmente na construção de \textbf{um} mundo mais \textbf{justo}, equitativo, democrático e humano.
O contexto escolar deve ser preparado visando a \textbf{uma} formação cidadã em \textbf{que} todas as crianças e adolescentes devem ser protegidos contra práticas \textbf{que} \textbf{fomentem} a exploração do trabalho infantil e discriminação étnico-racial, religiosa, sexual, de gênero, pessoa com deficiência ou de qualquer outra ordem.
O Processo de Envelhecimento, Respeito e Valorização do Idoso ( Lei nº 10.741/2003 ) é \textbf{um} tema \textbf{que} propõe entender o envelhecimento como \textbf{um} fenômeno natural da condição humana.
Para além da cronologia, há \textbf{um} conjunto amplo de aspectos \textbf{que} também \textbf{configuram} essa etapa do desenvolvimento humano : biológicos, culturais, históricos, psicológicos e sociais.
Embora o envelhecimento humano seja \textbf{uma} condição natural, as representações e sentimentos são construídos socialmente.
Dessa forma, faz -se necessário \textbf{que} as escolas incluam, em suas práticas curriculares, ações \textbf{que} visem ao desenvolvimento de comportamentos e atitudes \textbf{que} \textbf{aproximem} as gerações, estimulem os estudantes para o convívio, destituído de preconceitos, com pessoas idosas e sejam \textbf{educadas} para o envelhecimento humano.
O objetivo é garantir o respeito, a \textbf{dignidade} e a educação ao longo da vida.
Assim, no âmbito escolar, deve -se também reconhecer o protagonismo da pessoa idosa \textbf{enquanto} estudante e como sujeito \textbf{que}, munido de experiências e saberes, aprende mais sobre si mesmo e sobre o mundo.
A Educação Ambiental ( Lei nº 9.795/1999, Parecer CNE/CP nº 14/2012, Resolução CNE/CP nº 2/2012 e Programa de Educação Ambiental de Pernambuco-PEA/PE 2015 ) é \textbf{um} processo contínuo, dinâmico, participativo e interativo de aprendizagem das questões socioambientais.
Dessa forma, a Educação Ambiental constitui \textbf{uma} das dimensões do direito ao meio ambiente equilibrado e sustentável, prioridade na garantia da qualidade de vida das pessoas por meio de concepções e práticas inter/transdisciplinares, contínuas e permanentes, realizadas no contexto educativo.
Priorizando as questões ambientais, devemos \textbf{despertar} no estudante a importância de manter relações harmoniosas entre a sociedade e a natureza, preservando a biodiversidade e as culturas.
É nessa perspectiva \textbf{que} as atividades educativas devem envolver a escola e a comunidade em seu entorno, refletir sobre atitudes de proteção e preservação da natureza, dialogando por meio dos diferentes componentes curriculares.
A Educação para o Consumo e Educação Financeira e Fiscal ( Parecer CNE/CEB nº 11/2010 e Resolução CNE/CEB nº 7/2010 ) são temas \textbf{que} apontam para \textbf{abordagens} na escola, proporcionando ao estudante ter \textbf{uma} compreensão sobre finanças e economia, consumo responsável, processo de arrecadação financeira e a aplicação dos recursos recolhidos como também sua importância para o valor social dos tributos, procedência e destinação.
De modo geral, essas \textbf{abordagens} devem possibilitar ao estudante analisar, fazer considerações fundamentadas, tomar decisões e ter posições críticas sobre questões financeiras \textbf{que} envolvam a sua vida pessoal, familiar e da realidade social e, por \textbf{conseguinte}, compreender a cidadania, a participação social, a importância sobre as questões tributárias, o orçamento público, seu controle, sua execução e sua transparência, bem como a preservação do patrimônio público.
A Educação das Relações Étnico-raciais e Ensino da História e Cultura Afro-brasileira, Africana e Indígena ( Leis nº 10.639/2003 e 11.645/2008, Parecer CNE/CP nº 3/2004, Resolução CNE/CP nº 1/2004 e Parecer CNE/CEB nº 14/2015 ) é \textbf{uma} temática \textbf{que} deve ser trabalhada articulada a diferentes componentes curriculares, mas também no âmbito do currículo como \textbf{um} todo.
Deve assegurar o conhecimento e o reconhecimento desses povos na formação cultural, social, econômica e histórica da sociedade brasileira, ampliando as referências socioculturais da comunidade escolar na perspectiva da valorização da diversidade étnico-racial, contribuindo para a construção e afirmação de diferentes identidades.
É necessário \textbf{que} as práticas escolares contemplem nos seus currículos o ensino da história e cultura afro-brasileira, africanas e indígenas como forma de reconhecimento da contribuição \textbf{que} diversos povos deram para a história e cultura nacional.
Desta maneira, será alcançada \textbf{uma} educação das relações étnico-raciais \textbf{que} respeite a diversidade brasileira e \textbf{que} busque a erradicação da desigualdade e discriminação, ensejando a construção de \textbf{uma} sociedade \textbf{baseada} no reconhecimento das diferenças e na verdadeira democracia racial.
Ao abordarmos a Diversidade Cultural ( Parecer CNE/CEB nº 11/2010 e Resolução CNE/CEB nº 7/2010 ), biológica, étnico-racial, devemos considerar a construção das identidades, o contexto das desigualdades e dos conflitos sociais.
Este tema aborda a construção histórica, social, política e cultural das diferenças \textbf{que} estão ligadas às relações de poder, aos processos de colonização e dominação.
Este currículo propõe ações e práticas educativas \textbf{que} contemplem essa temática na sala de aula e em toda comunidade escolar para \textbf{que} se promova o combate ao preconceito e à discriminação.
É importante, no contexto escolar, possibilitar a compreensão de \textbf{que} a sociedade humana, sobretudo a brasileira, é composta por vários elementos \textbf{que} formam a diversidade cultural e a identidade de cada povo e de cada comunidade.
A partir dessa perspectiva, devem ser desenvolvidas atitudes de respeito às diferenças, considerando \textbf{que} a completude humana é construída na interação entre as diferentes identidades.
A Relação de Gênero ( Parecer CNE/CEB nº 07/2010, Resolução CNE/CEB nº 02/2012, Lei nº 11.340/2006 – Lei Maria da Penha, Plano Nacional de Educação em Direitos Humanos, 2006, Instrução Normativa da SEE nº 007/2017 e Portaria MEC nº 33/2018 ) é entendida como \textbf{uma} \textbf{categoria} de análise \textbf{que} ajuda a pensar a maneira como as ações e posturas dos \textbf{homens} e das mulheres são determinados pela cultura em \textbf{que} estão inseridos ( SCOTT, 1990 ).
Deve ser também compreendida como \textbf{um} conceito \textbf{baseado} em parâmetros científicos de produção de saberes \textbf{que} transversaliza diversas áreas do conhecimento, sendo capaz de identificar processos históricos e culturais \textbf{que} classificam e posicionam as pessoas a partir de \textbf{uma} relação sobre o \textbf{que} é entendido como feminino e masculino, essencial para o desenvolvimento de \textbf{um} olhar referente à reprodução de desigualdades no contexto escolar.
A perspectiva da ‘ igualdade de gênero ’, no currículo, é pauta para \textbf{um} sistema escolar inclusivo \textbf{que} crie ações específicas de combate às discriminações e \textbf{que} não contribua para a reprodução das desigualdades \textbf{que} persistem em nossa sociedade.
Não se trata, portanto, de anular as diferenças percebidas entre as pessoas, mas sim de fortalecer a democracia à medida \textbf{que} tais diferenças não se desdobrem em desigualdades.
A garantia desse debate e a elaboração de estratégias de enfrentamento às diversas formas de violência são, portanto, direitos assegurados por lei.
Esses são pautados em demandas emergenciais e \textbf{que} reafirmam a necessidade dos espaços escolares serem \textbf{lócus} de promoção da cidadania e respeito às diferenças.
Para efetivar isso, é necessária a implementação de ações com a perspectiva de eliminar atitudes ou comportamentos preconceituosos ou discriminatórios relacionados à ideia de inferioridade ou superioridade de qualquer orientação sexual, identidade ou expressão de gênero.
A Educação Alimentar e Nutricional ( Lei nº 11.947/2009 ) deve ser vivenciada por toda comunidade escolar de forma contínua e permanente, visando desenvolver práticas educativas, na perspectiva da segurança alimentar e nutricional, \textbf{que} respeitem a cultura, as tradições, os hábitos alimentares saudáveis e as \textbf{singularidades} dos estudantes.
Perpassa pela valorização da alimentação escolar, o equilíbrio entre qualidade e quantidade de alimentos consumidos, além do estudo sobre macro e micronutrientes necessários para a formação do indivíduo.
Dessa forma, o currículo traz a educação alimentar e nutricional, inserindo conceitos de alimentação e nutrição nas diferentes etapas de ensino, considerando o acesso à alimentação saudável como algo fundamental para o crescimento e desenvolvimento dos indivíduos.
Nessa dimensão, é necessário \textbf{que} o currículo desenvolva a \textbf{percepção} de \textbf{que} \textbf{uma} alimentação adequada e saudável é \textbf{um} direito humano, e \textbf{que} seja adquirida e consumida garantindo a segurança alimentar e nutricional.
A alta incidência de violência no trânsito, inclusive com mortes, \textbf{remete} à necessidade de incentivar a conscientização por meio de \textbf{um} trabalho de Educação para o Trânsito ( Lei nº 9.503/1997 ), envolvendo valores e princípios fundamentais para \textbf{um} convívio social saudável : respeito ao próximo, solidariedade, prudência e cumprimento às leis.
É \textbf{preciso} promover práticas educativas e intersetoriais \textbf{que} problematizem as condições da circulação e convivência nos espaços públicos desde a própria escola, seja no campo ou na cidade, para \textbf{que} se promova a convivência mais harmoniosa nos espaços compartilhados, de modo a incentivar \textbf{uma} circulação mais segura de forma eficiente e, sobretudo, mais humana.
A educação para o trânsito deve prever, no currículo da Educação Básica, a construção de valores direcionados ao comportamento respeitoso, ao cuidado com as pessoas e com o meio ambiente, considerando o direito humano à vida, \textbf{que} se constitui no seu bem maior.
Trazer o tema Trabalho, Ciência e Tecnologia ( Parecer CNE/CEB nº 11/2010 e Resolução CNE/CEB nº 7/2010 ) para o currículo da Educação Básica contribui para a compreensão do princípio educativo \textbf{que} envolve não só discussões acerca do mundo do trabalho, mas também acerca do desenvolvimento de capacidades humanas para transformação da realidade material, social.
Relaciona -se ainda à compreensão da Ciência e Tecnologia \textbf{enquanto} dimensões capazes de provocar reflexões e intervenções sobre o mundo nos aspectos sociais e naturais sem \textbf{perder} de vista o caráter da sustentabilidade.
Nesse sentido, é fundamental \textbf{que} os currículos e as práticas dos professores promovam a pesquisa, como princípio pedagógico, associada a \textbf{uma} \textbf{abordagem} reflexiva dos conteúdos \textbf{que} considere a relação complexa entre os potenciais do Trabalho, da Ciência e da Tecnologia para resolução de problemas, a ampliação da capacidade produtiva e empreendedora, bem como para a garantia de \textbf{um} espaço de reflexão e atuação crítica e ética sobre suas \textbf{influências} nos impactos ambientais e sociais.
É importante \textbf{que} o currículo da Educação Básica, ao abordar essa temática, promova \textbf{uma} reflexão sobre as diversas formas de trabalho, o uso das tecnologias, às suas respectivas funções e organização social em torno de cada profissão, a contribuição dessas para o desenvolvimento da sociedade, bem como sobre as relações sociais e de poder \textbf{que} se estabelecem em torno do mundo do trabalho.
A temática Saúde, Vida Familiar e Social ( Parecer CNE/CEB nº 11/2010, Resolução CNE/CEB nº 7/2010, Decreto nº 7.037/2009, Parecer CNE/CP nº 8/2012 e Resolução CNE/CP nº 1/2012 ) provoca a reflexão sobre conceitos \textbf{que} nos \textbf{remetem} não só a ausência de doença, mas, sobretudo, ao completo bem-estar experimentado por pessoas saudáveis.
A concepção \textbf{que} se entende por saúde tem relações diretas com o meio cultural, social, político, econômico, ambiental e afetivo em \textbf{que} se \textbf{vive}.
A visão histórica dos diversos significados de saúde também sofre variações ao longo do tempo.
O currículo, ao desenvolver esse tema, deve considerar a saúde numa perspectiva mais ampla \textbf{que} envolve as várias dimensões do ser humano, tais como : saúde mental, comportamental, atitudinal, orgânica, física, motora, afetiva, sensorial, entre outras.
É necessário \textbf{que} a pessoa se perceba em sua multidimensionalidade e \textbf{que} a esfera da saúde seja reconhecida sob os diversos aspectos \textbf{que} envolvem \textbf{uma} vida saudável.
O contexto político relativo a como a sociedade está organizada também interfere na dimensão da saúde do cidadão.
A estrutura da saúde pública, o planejamento das cidades, o saneamento básico, o estilo de vida do/no campo ou da/na cidade, o sistema de transporte e habitacional, as relações familiares e sociais poderão interferir na saúde das pessoas.
Esses aspectos devem ser considerados e refletidos no currículo de forma a levar os estudantes a compreenderem e buscarem \textbf{um} estilo de vida mais saudável.
Assim sendo, os temas integradores, acima abordados, além de estarem presentes em habilidades e competências de diferentes componentes curriculares, devem estimular o desenvolvimento de atividades para serem vivenciadas no contexto da escola, envolvendo todas as áreas do conhecimento \textbf{que} compõem o currículo.
Por isso, é necessário \textbf{que} se realize \textbf{um} trabalho interdisciplinar, motivador, inclusivo, resultando em \textbf{uma} experiência mais enriquecedora para os estudantes, os professores participantes e também toda a comunidade escolar.

% matched lemmas: abordagem, aproximar, basear, categoria, compromisso, configurar, conseguinte, conseguir, demandar, despertar, dignidade, direção, educar, embasar, enquanto, fomentar, homem, influência, justo, lócus, norteador, percepção, perder, preciso, qualificar, que, remeter, singularidade, sujeito, um, ver, viver
\end{textsample}
